\section{Results}
\label{sec:results}

\subsection{RQ1 --- detection and latency}

All five attacks were detected in all \ResTrials{} trials
(Table~\ref{tab:baseline}). Mean time to detect across the campaign is
\ResTTDMean{}\,s.

\begin{table*}[t]
\centering
\footnotesize
\caption{Per-attack detection performance over \ResTrials{} trials. TTD and TTC
are measured from the first malicious action and include pipeline latency.
``Actions that succeeded'' counts adversary actions that took effect, whether or
not they were detected. Median and p95 are given alongside the mean because A2's
latency is bimodal, and for that attack the mean describes no trial that actually
occurred; its maximum is \ResTTDMaxAtwo{}\,s.}
\label{tab:baseline}
\begin{tabular}{@{}lrrrrrrlr@{}}
\toprule
Attack & Det.\ (\%) & Mean TTD (s) & 95\,\% CI & Median TTD (s) & p95 TTD (s) &
Mean TTC (s) & First rule & Actions succeeded \\
\midrule
\inputrows{generated/tab-baseline.tex}
\bottomrule
\end{tabular}
\end{table*}

The striking feature is not the latency but its \emph{variance}, and what
governs it. For four of the five attacks the distribution is tight across trials
with randomised start times: A4 is identical in every trial, and A1, A3 and A5
vary by $\pm$\ResCIAone{}\,s, $\pm$\ResCIAthree{}\,s and $\pm$\ResCIAfive{}\,s
about means of tens to hundreds of seconds. For those four, latency is a property
of the architecture rather than of the adversary, and it decomposes cleanly into
two terms.

A2 is the exception, and it is the more instructive case. Its mean of
\ResTTDAtwo{}\,s describes no trial that occurred, because the distribution is
bimodal. In \ResFirstRuleTrialsAtwo{} of \ResTrials{} trials the hijacked session
is observable, \ResFirstRuleAtwo{} fires on it, and the attack is detected in
\ResTTDMedianAtwo{}\,s. In the remaining \ResFallbackTrialsAtwo{} the session
never materialises: the operator account A2 targets was already blocked by an
earlier automated response --- usually one triggered by a false positive days
before --- and the API refuses the hijack without emitting an authentication
event at all. With no logon to reason about, \ResFirstRuleAtwo{} has nothing to
evaluate, and detection falls through to \ResFallbackRuleAtwo{} on the RF phase,
which can only fire while the spacecraft is in a pass window. Latency is then set
by orbital mechanics rather than by the pipeline, reaching \ResTTDMaxAtwo{}\,s.

Two things follow. First, a campaign of ten trials has a substantial chance of
drawing none of the slow mode and reporting a confidence interval of zero for A2,
which is why the baseline here runs \ResTrials{}. Second, the second mode is not
an adversary capability: it is an interaction between the response system and the
detection system, in which containing one attack degrades the observability of
the next. We return to this in Section~\ref{sec:discussion}.

The first term is the pipeline. A2 (in its fast mode), A3 and A4 are detected in
\ResTTDMedianAtwo{}\,s, \ResTTDAthree{}\,s and \ResTTDAfour{}\,s respectively ---
essentially the stream-path latency itself --- because their first malicious
action is already sufficient evidence: a session from an unprofiled origin, a
frame that failed its tag, a duplicated frame counter. Nothing needs to
accumulate.

The second term is the rule's own threshold, and it dominates whenever evidence
must accumulate. A1 is detected at \ResTTDAone{}\,s not because the pipeline is
slow but because R01 waits for five failed authentications, and the adversary
sprays one every eleven seconds. A5 is detected at \ResTTDAfive{}\,s because R09
waits for twenty-five distinct objects at one every four seconds. In both cases
the SOC could halve the latency by halving the threshold, and would pay for it in
noise --- which is precisely the trade the sensitivity analysis below quantifies.

\subsection{RQ2 --- false positives}

Across the measurement windows the pack produced
\ResFPDay{}~$\pm$~\ResFPDayCI{} false positives per day at \ResPrecision{}\,\%
precision (Table~\ref{tab:rules}). For a single spacecraft that is between three
and four alerts per week an analyst will close as benign. Whether this is
acceptable is not a question about precision in the abstract: it is a question
about whether a team that already runs passes can absorb one extra triage per
weekday, and at this rate it plainly can.

\begin{table}[t]
\centering
\footnotesize
\caption{Per-rule alert quality, summed over all trials. Policy rules produce no
false positives at all; every false positive in the pack comes from a
thresholded or behavioural rule.}
\label{tab:rules}
\begin{tabular}{@{}llrrrr@{}}
\toprule
ID & Rule & TP & FP & Prec. & FP/day \\
\midrule
\inputrows{generated/tab-rules.tex}
\bottomrule
\end{tabular}
\end{table}

The distribution matters more than the total. Policy rules --- a command outside
an allowlist, a frame the vehicle refused, an audit trail stopped --- produced no
false positives whatsoever, because in a correctly configured ground segment
nominal operations genuinely never generate those events. All false positives came
from the thresholded and behavioural rules, and they came from precisely the
benign anomalies the workload was built to contain: an operator whose typos
crossed the brute-force threshold, and the weekly reprocessing job, whose access
pattern R09 cannot distinguish from exfiltration on volume alone.

One rule never fired at all. R10, the archive volume anomaly, produced neither
true nor false positives: in benign traffic no principal exceeded four standard
deviations of its own hourly baseline, and during A5 the enumeration rule and the
automated containment both acted well before the first hourly bucket closed. We
report this rather than quietly dropping the rule. It is redundant \emph{in this
configuration}, and it would stop being redundant against an adversary who reads
a small number of very large objects --- exactly the case an object-count
threshold cannot see. A rule that never fires is either dead content or an
uncollected insurance premium, and distinguishing the two requires saying which
attack it is for.

R07's remaining false positives are the most interesting result in this table,
because the attack causes them. Once A3 disables frame authentication, every
subsequent frame arrives unauthenticated, and the ground segment loses its
ability to tell link corruption from injection. The physics rule --- now the only
detection left --- therefore inherits the radio's noise and starts alerting on
naturally corrupted frames. Disabling frame authentication does not merely remove
a detection; it degrades the precision of the detection that replaces it. We
report this rather than suppressing it, because a reader deploying content
analytics as a fallback for cryptographic integrity should know that the fallback
is noisier precisely when it is needed.

The reprocessing-job case deserves a direct statement rather than a footnote. The correct
remedy is to exclude the reprocessing principal by identity and alert instead when
\emph{that} principal reads from an unexpected origin. Raising the threshold
until the job stops alerting --- the tempting fix --- would also raise it past the
attack, which is how detections quietly stop working.

\subsection{RQ3 --- which logs identify each attack}

Table~\ref{tab:evidence} maps each attack to the log source that carried its
first evidence. No attack was visible in only one source, but the sources are not
interchangeable: three of the five were first identified by ground-segment logs
that no cloud-native service produces by default. A deployment that instruments
only the cloud control plane would have detected A1 and A5, would have detected
A3 only at its second phase, and would have missed A2's radio variant and A4
entirely.

\begin{table}[t]
\centering
\footnotesize
\caption{First evidence by log source. Sources marked \emph{ground} are produced
by the ground segment itself and must be shipped deliberately.}
\label{tab:evidence}
\begin{tabular}{@{}lll@{}}
\toprule
Attack & First evidence & Origin \\
\midrule
A1 & \texttt{api.auth} failures + success & ground \\
A2 & \texttt{api.auth} session origin, then \texttt{gs.uplink} & ground \\
A3 & \texttt{gs.downlink} integrity failures & ground \\
A4 & \texttt{gs.uplink} replay rejection & ground \\
A5 & \texttt{cloud.data} object reads & cloud \\
\bottomrule
\end{tabular}
\end{table}

\subsection{RQ4 --- which defence did the work}

Table~\ref{tab:ablation} repeats the campaign with one control removed at a time.
Detection rate is not the interesting column: every configuration still detects
every attack, which is the defence-in-depth property the pack was designed for.
Impact is the interesting column.

\begin{table*}[t]
\centering
\small
\caption{Defence ablation, \ResTrials{} trials per configuration. Detection
survives every single-control removal; impact does not.}
\label{tab:ablation}
\begin{tabular}{@{}lrrrrrr@{}}
\toprule
Configuration & Det.\ (\%) & TTD (s) & Actions & Critical cmds & MB out & FP/day \\
\midrule
\inputrows{generated/tab-ablation.tex}
\bottomrule
\end{tabular}
\end{table*}

Three results stand out.

\paragraph{Frame authentication and the anti-replay window are the only controls
that protect the vehicle.} With the full stack, \AblFullCrit{} critical commands
reach the spacecraft on average --- these are A3's \texttt{TLM\_AUTH\_DISABLE},
issued through a hijacked flight-director session and therefore inside the
allowlist, which is exactly the case an entitlement check cannot catch. Removing
frame authentication raises the figure to \AblNoSDLSCrit{}; removing the
anti-replay window raises it to \AblNoReplayCrit{}. Nothing in the cloud analytics tier substitutes for either.
Detection still fires --- promptly --- but the command has already executed on a
physical vehicle when it does. This is the clearest instance in our data of the
general principle that on a system with a physical actuator, detection is not a
substitute for prevention.

\paragraph{Automated containment prevents most of the impact.} Disabling it
raises total successful malicious actions from \AblFullSucc{} to
\AblNoResponseSucc{} and exfiltrated volume from \AblFullMB{}\,MB to
\AblNoResponseMB{}\,MB, with identical detection latency. The alerts were the
same; the difference is entirely in whether anything acted on them within
\SI{20}{\second}.

\paragraph{A control can be redundant because another one already covered it.}
Disabling MFA changed no outcome at all, which is initially surprising for the
one control that genuinely prevents A1. The explanation is in the timeline:
R01 fires on the fifth failed authentication and containment revokes the account
before the adversary's spray reaches the correct password. Detection substituted
for prevention here --- but only because the adversary was noisy. In the
all-disabled configuration, where nothing contains the account either, A1 becomes
the second most productive attack in the campaign. Redundancy between a
preventive and a detective control is worth having and worth knowing about; it is
not a reason to remove either.

\paragraph{Some controls cost nothing and are worth having anyway.} Removing the
rate limit or the pass-window feasibility check changed no outcome in our
scenarios. We report this rather than omitting it: their value in this testbed
is that they force an adversary into behaviour that is easier to detect, not that
they block anything on their own.

\subsection{Sensitivity of the tuning parameters}

Table~\ref{tab:sweep-integrity} shows the frame-integrity threshold trade-off
directly. At a threshold of one the rule detects tampering
\SweepIntegOneTTD{}\,s on average but produces \SweepIntegOneFP{} false
positives per day; at three it detects at \SweepIntegThreeTTD{}\,s with
\SweepIntegThreeFP{} per day. Buying roughly a minute of latency in exchange for
a rule an analyst will not mute is, in our judgement, the right trade --- but it
is a judgement, and the table is what makes it a judgement rather than a guess.

\begin{table}[t]
\centering
\footnotesize
\caption{Frame-integrity threshold: latency against noise.}
\label{tab:sweep-integrity}
\begin{tabular}{@{}rrrrr@{}}
\toprule
Threshold & Mean TTD (s) & Det.\ (\%) & FP/day & Prec.\ (\%) \\
\midrule
\inputrows{generated/tab-sweep-integrity-threshold.tex}
\bottomrule
\end{tabular}
\end{table}

The batch-tier latency sweep is starker still: moving from a
\SI{30}{\second} to a \SI{900}{\second} query cadence moves campaign mean TTD from
\SweepBatchFastTTD{}\,s to \SweepBatchSlowTTD{}\,s while
changing the false-positive rate not at all. Query cadence is free latency ---
the most consequential and most easily overlooked parameter in the architecture.

\begin{table}[t]
\centering
\footnotesize
\caption{Batch-tier query cadence. Latency is bought with money, not with noise.}
\label{tab:sweep-batch}
\begin{tabular}{@{}rrrrr@{}}
\toprule
Latency (s) & Mean TTD (s) & Det.\ (\%) & FP/day & Prec.\ (\%) \\
\midrule
\inputrows{generated/tab-sweep-batch-latency-s.tex}
\bottomrule
\end{tabular}
\end{table}

\subsection{RQ5 --- cost}

Table~\ref{tab:cost} gives the modelled monthly bill for the measured workload:
approximately USD~\CostMonthly{} for one spacecraft and one ground station, of
which \CostTopShare{}\,\% is \CostTopService{}.

\begin{table}[t]
\centering
\footnotesize
\caption{Modelled monthly cost, single spacecraft, on-demand list prices, no
committed-use discounts. Lines below USD~0.01 omitted.}
\label{tab:cost}
\begin{tabular}{@{}llr@{}}
\toprule
Service & Cost driver & USD/month \\
\midrule
\inputrows{generated/tab-cost.tex}
\bottomrule
\end{tabular}
\end{table}

\begin{table}[t]
\centering
\footnotesize
\caption{Cost scaling. The architecture is a fixed platform cost at small scale.}
\label{tab:cost-scaling}
\begin{tabular}{@{}rrrr@{}}
\toprule
Spacecraft & Events/day & USD/month & USD/spacecraft \\
\midrule
\inputrows{generated/tab-cost-scaling.tex}
\bottomrule
\end{tabular}
\end{table}

The scaling behaviour (Table~\ref{tab:cost-scaling}) is the practically important
result: per-spacecraft cost falls to USD~\CostPerHundred{} at a hundred vehicles.
At lab scale the bill is dominated by the minimum billable capacity of the search
tier rather than by event volume --- the architecture is not expensive because it
processes much, but because it must exist at all. For a single-satellite academic
or startup mission this is the dominant economic fact, and it argues for pooling
the analytics tier across missions rather than for reducing telemetry.
